\subsection{Rank bounds from the exponent}\label{sec:finite-ranks}
We use the passage from an algorithm at one fixed size to an algorithm
on tensor powers, together with the closure argument for border rank
\cite{bini,bcs,spectrum}.

\begin{proof}[Proof of Corollary~\ref{cor:tight}]
Fix \(n\ge2\). A rank-\(r\) decomposition of \(T_{n^j}\) over any field
\(K\) gives \(\omega_K\le\log_{n^j} r\) by the recursion and padding in
Lemma~\ref{lem:specialize}. Theorem~\ref{thm:square} therefore gives
\(\operatorname{rank}_K(T_{n^j})\ge n^{9j/4}\) at each \(j\). Taking
roots gives the asymptotic-rank lower bound. At \(K=\overline F\) and
\(j=1\), Lemma~\ref{lem:bilinearize} now gives the claimed lower bound
of \(\lceil n^{9/4}/2\rceil\) on the number of nonlinear gates.

For the geometric border bound, work in the fixed tensor space \(V\)
over \(\overline F\) containing \(T_n\). Suppose \(T_n\) is in the
Zariski closure of the rank-at-most-\(r\) tensors. For each \(j\),
their \(j\)-th powers span a space containing \(T_n^{\otimes j}\),
since membership of a tensor power in this span is a closed condition.
The number of coefficients of a homogeneous polynomial of degree \(j\)
in \(\dim V\) variables bounds the dimension of the span. We can thus
express \(T_n^{\otimes j}\) using \(O_n(j^{\dim V-1})\) such powers,
each of rank at most \(r^j\). Comparison with the lower bound at
\(n^j\) and passage to \(j\)-th roots give \(r\ge n^{9/4}\).

For the asymptotic-rank upper bound fix \(\varepsilon>0\). At side \(n^j\)
the known circuit upper bound and Lemma~\ref{lem:bilinearize} give a
decomposition of rank \(r_j=O_{F,\varepsilon}(n^{j(9/4+\varepsilon)})\)
over \(\overline F\). Its coefficients belong to a finite extension
\(E_j/F\), of degree \(d_j\). For each \(t\), power this decomposition
\(t\) times. Expand each of the three vectors in each simple tensor
in an \(F\)-basis of \(E_j\) containing \(1\), then project the
products of basis coefficients onto the coordinate of \(1\).
The tensor being represented is defined over \(F\), so this gives
\[
        \operatorname{rank}_F(T_n^{\otimes jt})\le d_j^3 r_j^t .
\]
The limit defining asymptotic rank exists by submultiplicativity.
First take \(jt\)-th roots and let \(t\to\infty\) with \(j\) fixed.
Then let \(j\to\infty\) and \(\varepsilon\downarrow0\). The resulting
upper bound is \(n^{9/4}\) over the original field.
\end{proof}

At finite sizes, Corollary~\ref{cor:tight} also conflicts with the dyadic
border-rank estimate proposed in the first version of Sun's
preprint~\cite{Sun}. At $n=2^8$ the proposed upper bound is $2^{17}$
and our lower bound is $2^{18}$.

The square tensor also decides a prediction about tight three-tensors.
A support is tight when each local basis admits distinct integer
weights whose sum is zero on every support triple. In \(T_n\), take
\(B>n\) and assign weights \(i-Bj\), \(Bj-B^2k\) and \(B^2k-i\)
to \(e_{ij},e_{jk},e_{ki}\), respectively. Each assignment is
injective on its factor and the three weights sum to zero. Each
one-factor flattening has rank \(n^2\). Strassen's conjecture for
tight tensors predicts asymptotic rank \(n^2\) in this case
\cite[Conjecture~5.3]{Strassen1994}. Corollary~\ref{cor:tight}
gives \(n^{9/4}\).

\begin{proof}[Proof of Corollary~\ref{cor:rectangular}]
The three one-factor flattenings have ranks \(ab,bc,ca\).
These are lower bounds for geometric border rank by the closedness
of matrix-rank conditions, and for asymptotic rank by multiplicativity
of flattening ranks. Grouping corresponding factors in
\[
            T_{a,b,c}\otimes T_{b,c,a}\otimes T_{c,a,b}
                                \simeq T_{abc}
\]
turns each of the three indices into a tuple with \(abc\) possibilities.
Both ranks are invariant under cyclic relabeling of the factors and
submultiplicative under tensor product. Hence each of their values
on \(T_{a,b,c}\), cubed, is at least \((abc)^{9/4}\) by
Corollary~\ref{cor:tight}. The case \(abc=1\) is immediate.

For \(a=c=n\) and \(b=\lfloor n^\beta\rfloor\), the geometric border
bound is \(\Omega(n^{3(2+\beta)/4})\). The bilinearization in
Lemma~\ref{lem:bilinearize} applies entrywise to rectangular products
as well. A quadratic exponent therefore requires \(3(2+\beta)/4\le2\).
This gives \(\alpha_F\le2/3\), and the upper bound for multiplication
at \(\beta=1/2\) gives \(\alpha_F\ge1/2\)~\cite{abvz}.
\end{proof}
